\chapter{Considerações Finais}
\label{cap:conclusao}

Este trabalho investigou empiricamente a aplicação do método USARP na elicitação de requisitos de usabilidade, por meio de uma síntese empírica de braço único que integrou evidências de quatro estudos primários conduzidos em contextos acadêmicos e industriais. A motivação partiu de uma lacuna concreta: embora estudos isolados indicassem benefícios do método, não havia uma síntese que avaliasse sua aplicação em diferentes contextos e identificasse padrões recorrentes de percepção, uso e adoção. Para preencher essa lacuna, consolidaram-se, no nível do respondente, os dados de levantamento dos estudos que adotaram o instrumento comum, e essa base foi submetida a uma reanálise quantitativa, com estatística descritiva, redução de dimensionalidade e testes não-paramétricos, e qualitativa, com análise temática, articulada com as evidências reportadas por \citeonline{oliveirajr2020}, \citeonline{marques2022}, \citeonline{fiori2022} e \citeonline{simoes2022usarp}.

No plano quantitativo, a percepção sobre o USARP mostrou-se elevada e marcadamente homogênea entre contextos e coortes, com forte efeito de teto em todos os dez itens analisados. A redução de dimensionalidade revelou uma primeira dimensão de percepção positiva geral, responsável por cerca de um terço da variância, acompanhada de contrastes secundários entre os blocos de cartas (requisitos, mecanismos e prototipação) e os itens de percepção operacional. Nem a segmentação não supervisionada, que agrupou os respondentes por nível geral de avaliação e não por contexto, nem as comparações inferenciais, que não acusaram diferenças robustas após a correção de Holm, sustentaram uma distinção marcante entre os ambientes acadêmico e industrial. A única diferença que resistiu ao controle estatístico ocorreu no item de contexto do mecanismo de usabilidade, entre a coorte industrial de 2022 e duas turmas acadêmicas.

No plano qualitativo, a análise temática das respostas abertas confirmou as cartas, o checklist e os guias de especificação como o principal motor da elicitação, seguidos pela organização e pela padronização do processo, pela colaboração entre profissionais de diferentes papéis e pelo uso conjunto de personas e \textit{user stories}. As tensões concentraram-se na curva de aprendizado, no volume de cartas e na linguagem técnica de UI/UX, e as adaptações sugeridas foram majoritariamente incrementais, sem propostas de revisão estrutural do método. Um achado convergente entre os estudos e entre as duas naturezas de evidência foi a subutilização das cartas de prototipação, menos citadas nos relatos e com a menor média entre os itens de cartas.

Esses resultados atendem aos objetivos específicos propostos: sintetizaram-se as evidências de percepção de utilidade, produtividade e eficácia (primeiro objetivo); caracterizou-se o uso efetivo dos mecanismos, com destaque para a centralidade das cartas e a subutilização da prototipação (segundo objetivo); evidenciaram-se mais semelhanças do que diferenças entre contextos (terceiro objetivo); extraíram-se as estruturas latentes associadas à aplicação do método por meio da Análise de Componentes Principais (quarto objetivo); e discutiram-se as implicações práticas e as limitações (quinto objetivo). Atende-se, assim, ao objetivo geral de investigar empiricamente a aplicação do USARP a partir da síntese de evidências de contextos acadêmicos e industriais.

Como contribuições, o trabalho oferece uma síntese empírica crítica do USARP, que reúne, de forma comparável e no nível do respondente, evidências antes dispersas em estudos individuais; consolida e documenta o instrumento de avaliação derivado do \textit{Technology Acceptance Model} e da utilidade percebida das cartas; disponibiliza um \textit{pipeline} de análise reprodutível, com critérios e diagnósticos explícitos; e fornece subsídios práticos para a evolução do método, ao identificar tanto seus pontos fortes quanto os elementos que ainda impõem maior carga de uso.

O estudo apresenta limitações que devem ser consideradas na interpretação de seus resultados. A amostra é de tamanho reduzido, e a subamostra industrial, em particular, é pequena, com onze respondentes, o que limita o poder estatístico das comparações entre contextos; nesse sentido, a ausência de diferenças significativas indica ausência de evidência de diferença, e não equivalência comprovada. O forte efeito de teto e o valor modesto da medida de Kaiser-Meyer-Olkin recomendaram o uso exploratório, e não confirmatório, da Análise de Componentes Principais. Além disso, por se tratar de uma síntese de braço único baseada em dados secundários, coletados por instrumentos não idênticos entre todos os estudos e sem grupo de comparação controlado, não se formulam afirmações causais, e a generalização permanece circunscrita a cenários semelhantes aos estudados. Essas e outras ameaças à validade foram detalhadas no Capítulo \ref{cap:metodologia}.

Como trabalhos futuros, recomenda-se a condução de estudos observacionais sobre o uso efetivo do método, capazes de complementar as medidas de percepção com evidências de comportamento, conforme já sugerido por \citeonline{marques2022}. A ampliação da base industrial, com coortes maiores, permitiria comparações entre contextos com maior poder estatístico. No plano do próprio método, os achados apontam direções concretas de refinamento: reduzir a carga associada ao volume de cartas, por exemplo por meio de uma etapa de filtragem assistida e de descrições contextuais, na forma de tooltips, no checklist; revisitar as cartas de prototipação, consistentemente subutilizadas; e investigar em que medida o método contribui para a prevenção de problemas de usabilidade e para a redução de retrabalho ao longo do ciclo de desenvolvimento. Estudos longitudinais poderiam, ainda, examinar como a maturidade de uso afeta a percepção dos mecanismos, dimensão em que se observou a maior variação entre as coortes.

Em síntese, as evidências reunidas indicam que o USARP é percebido de maneira consistentemente positiva e similar em contextos acadêmicos e industriais, tendo nas cartas e nos guias o seu principal valor percebido, e na curva de aprendizado e no volume de artefatos os seus principais desafios. Ao consolidar e interpretar essas evidências de forma integrada, espera-se que este trabalho contribua para a adoção informada do método e para a continuidade de sua evolução, em benefício da incorporação sistemática da usabilidade desde as fases iniciais do desenvolvimento de software.